% !TeX program = pdflatex
% ─────────────────────────────────────────────────────────────────────────────
%  papers/cripto.tex — a Cadeia de Cristais (CristalChain):
%  livro-razão cuja cifra conserva
%
%  Nome formal: Cadeia de Cristais. Marca / repo / catálogo: CristalChain.
%  Resgata o enunciado já construído em catalogo.tex («A CristalChain ---
%  um livro-razão cuja cifra conserva»). Não é blockchain. Não é PoW.
%  Não é palco quântico. Não é U. Não é 15.ª Parte. Sem Lei 8.
%
%  Compila:  cd papers && pdflatex cripto.tex
% ─────────────────────────────────────────────────────────────────────────────
\documentclass[11pt,a4paper]{article}
\usepackage[utf8]{inputenc}\usepackage[T1]{fontenc}\usepackage[brazil]{babel}
\usepackage{amsmath,amssymb,amsthm}
\usepackage[margin=2.4cm]{geometry}
\usepackage{booktabs}\usepackage{array}\usepackage{xcolor}
\usepackage[htt]{hyphenat}
\usepackage[hidelinks,breaklinks]{hyperref}

\sloppy
\emergencystretch=2em

\definecolor{cinza}{gray}{0.35}
\newcommand{\code}[1]{\texttt{\small #1}}
\newcommand{\caixa}[1]{%
  \par\smallskip\noindent
  \begin{center}
  \fbox{\begin{minipage}{0.94\textwidth}\centering\smallskip #1\smallskip\end{minipage}}
  \end{center}\smallskip}
\newcommand{\abs}[1]{\lvert #1\rvert}
\newcommand{\medido}[1]{\par\medskip\noindent{\small\color{cinza}\textbf{Medido.} #1}\par\medskip}
\newcommand{\fis}[2]{\textbf{#2}~\textnormal{(\code{#1})}}
\providecommand{\pt}[1]{\textbf{P#1}}
\newcommand{\sepcol}{\quad\penalty0$\big|$\quad\penalty0}
\newcommand{\colunas}[3]{%
\par\smallskip\noindent{\footnotesize\raggedright
\textbf{prova}: #1\sepcol\textbf{realiza}: #2\sepcol\textbf{verifica}: #3\par}\smallskip}
\providecommand{\interfacen}[1]{}
\newtheorem{teorema}{Teorema}
\newtheorem{lema}[teorema]{Lema}
\newtheorem{definicao}[teorema]{Definição}
\newtheorem{corolario}[teorema]{Corolário}
\newtheorem{proposicao}[teorema]{Proposição}
\newtheorem{observacao}[teorema]{Observação}
\renewcommand{\proofname}{Demonstração}

% ─────────────────────────────────────────────────────────────────────────────
%  papers/gkcapa.tex — a capa do Reino, mínima, para os papers em `article`.
%
%  O estilo.tex completo é do `report` (títulos de capítulo, skak, …). Aqui fica
%  só o que a capa precisa, para o paper continuar a compilar sozinho com
%  pdflatex. O tradutor próprio (tests/tex) carrega o estilo.tex da raiz / o
%  symlink papers/estilo.tex e avalia o mesmo `\gkcapa`.
% ─────────────────────────────────────────────────────────────────────────────
\definecolor{ouro}{HTML}{B8912F}
\definecolor{ouroclaro}{HTML}{F3E9CE}
\definecolor{tinta}{HTML}{1E1B16}
\definecolor{regua}{HTML}{6E6858}
\providecommand{\gknota}{\fontsize{7.62}{11.05}\selectfont}
\providecommand{\gktexto}{\fontsize{10.50}{15.22}\selectfont}
\providecommand{\gksub}{\fontsize{12.33}{17.87}\selectfont}
\providecommand{\gksec}{\fontsize{14.47}{20.98}\selectfont}
\providecommand{\gkcap}{\fontsize{16.99}{24.63}\selectfont}
\providecommand{\gktit}{\fontsize{23.42}{33.95}\selectfont}
\providecommand{\Prj}{\mathbb{P}}
\providecommand{\Trf}{\mathcal{T}}
\providecommand{\gkcapa}[3]{%
\title{\vspace{-0.6cm}
{\color{ouro}\rule{\textwidth}{1.4pt}}\\[6mm]
\textbf{\gktit\color{tinta} Reino Dourado}\\[3mm]{\gksec\itshape\color{regua} Golden Kingdom}\\[8mm]
{\gkcap\scshape\color{ouro} #1}\\[5mm]
{\gksub\color{regua} #2}\\[2mm]
{\gktexto\color{regua}\itshape #3}\\[7mm]
{\color{ouro}\rule{\textwidth}{0.6pt}}\\[9mm]{\gknota\color{regua}\begin{minipage}{\textwidth}\setlength{\parindent}{0pt}%
\textbf{Aviso de Isenção Algorítmica:} O universo, as leis e as entidades contidas nestas páginas ---
incluindo felinos matriciais, aves de rapina algébricas e amuletos de controle de fluxo --- são
construções de pura ficção formal. Esta obra não descreve a realidade física, não serve como
engenharia reversa do cosmos e não constitui doutrina religiosa. Qualquer semelhança com bugs reais do seu
sistema operacional é mera coincidência (ou falta de reza). Prossiga por sua própria conta e
risco.\end{minipage}}}
\author{}\date{}}


\begin{document}
\interfacen{6}
\gkcapa{Criptografia --- a Cadeia de Cristais (CristalChain)}
       {livro-razão cuja cifra conserva}
       {cripto $\neq$ palco quântico $\neq\mathcal{U}$; sem 15.\textsuperscript{a} Parte}
\maketitle

\begin{abstract}\noindent
A \textbf{Cadeia de Cristais (CristalChain)} \textbf{já estava construída}
no catálogo: um livro-razão cuja cifra conserva. Este paper \textbf{não}
inventa a cadeia, \textbf{não} a promove a blockchain, \textbf{não} a
funde com Born, Duo, Lei~8 nem com $A_0$ dourada. Resgata o enunciado,
fecha os labels (\code{cripto:}\ldots) e marca \texttt{nao localizada}
o elo que o repositório não mede.
\[
\boxed{\;
\pi^{2}=\pi
\neq x^{2}=1
\neq x^{2}=x+1
\neq x^{2}=-1
\neq\operatorname{Star}(\mathcal{U})=\mathcal{D}
\;}
\]
A conservação da cadeia é $\det(B^{k})=+1$ no ouro branco
(\code{tests/ouro\_branco.c}) --- o sinal do determinante sob
composição, \textbf{não} um saldo económico.
$P=k/N$ é GLH (\code{fis:thm:central}).
Born é $p_i=\lvert\langle i\rvert\psi\rangle\rvert^{2}$.
Parseval do corpus cristal é isometria do jornal
(\code{tests/cristal\_energia.js}).
$\mathrm{sha256}$ selecciona identidade
(\code{univ:def:lei-local-canonica}), \textbf{não} é a cifra.
Quatro conservações; não se fundem.
O fecho de $\mathcal{U}$ já está em
\code{univ:thm:composicao-fechos}; Born não o fecha.
\end{abstract}

% ═══════════════════════════════════════════════════════════════════════════
\section{I0 --- três palcos, cinco quadráticas}\label{cripto:i0}

\begin{observacao}[cripto $\neq$ palco quântico $\neq\mathcal{U}$]\label{cripto:obs:i0}
Este documento é \textbf{realização}: consome o construído
(\code{catalogo.tex}; \code{tests/ouro\_branco.c};
\code{tests/smartcontract.c}).
Não é Parte nova. Não reabre I0 das Partes fechadas.
Não inventa Lei~8. Não contradiz
\code{univ:thm:composicao-fechos}.
\caixa{%
criptografia (este paper)
$\;\neq\;$
palco quântico (\code{fis:obs:quantica-grafo-fecho})
$\;\neq\;$
$\mathcal{U}=(X,X^{*},\mathsf{Mor},\mathsf{Aut})$
(\code{univ:def:U})%
}
\end{observacao}

\begin{observacao}[cinco quadráticas, não se fundem]\label{cripto:obs:cinco-quadraticas}
\caixa{%
\begin{tabular}{@{}lll@{}}
$\pi^{2}=\pi$ & Alonzo / idempotente & \code{univ:def:alonzo-idemp}\\
$x^{2}=1$ & involução $A_0$ & \textbf{não} Duo\\
$x^{2}=x+1$ & $\operatorname{Star}(K)$ / ouro & slot I9; \textbf{não} $\operatorname{Star}(\mathcal{U})$\\
$x^{2}=-1$ & Lei~5 & \textbf{não} a estrela\\
$\operatorname{Star}(\mathcal{U})=\mathcal{D}$ & estrutural, bit $a=1$ & \textbf{não} elemento\\
\end{tabular}%
}
A família metálica é $x^{2}=nx+1$ e inclui o ouro $x^{2}=x+1$.
O ouro branco é $x^{2}=4x-1$ --- \textbf{outro lado} do determinante,
não mais um $n$ da família (\code{tests/ouro\_branco.c} §B1).
$A_0$ ($x^{2}=1$) $\neq\mathbf{Duo}$.
O associador do palco quântico
(\code{hiper/circular/associador\_ibm\_n2.py})
\textbf{não} é $A_0$ dourada.
\end{observacao}

\begin{observacao}[quatro conservações, não se fundem]\label{cripto:obs:conservacoes}
\caixa{%
\begin{tabular}{@{}lll@{}}
$\det=+1$ & cifra ($\det(B^{k})$; \textbf{não} saldo) & ouro branco; §\ref{cripto:cifra}\\
$P=k/N$ & rectângulo GLH & \code{fis:lem:conserva}; \code{fis:thm:central}\\
$\sum X\cdot X_{-}=N\sum x^{2}$ & Parseval do corpus & \code{tests/cristal\_energia.js}\\
$p_i=\lvert\langle i\rvert\psi\rangle\rvert^{2}$ & Born & \textbf{não} fecha $\mathcal{U}$\\
\end{tabular}%
}
$P=k/N\neq\lvert\psi_i\rvert^{2}$.
Born isolado:
\code{univ:thm:composicao-fechos};
\code{fis:obs:quantica-grafo-fecho}.
$\lvert\det\rvert=1$ da Lei~4 / volume
(\code{fis:thm:gato}(2))
é o invariante de $\operatorname{Aut}(T^{2})$,
\textbf{não} o $\det=+1$ da matriz $2\times2$ do ouro branco
--- homónimo; I0 decide.
Nenhuma das quatro linhas é conservação de valor económico.
\end{observacao}

% ═══════════════════════════════════════════════════════════════════════════
\section{Definição --- o que a cadeia já era}\label{cripto:def}

O catálogo já o escreveu
(\code{catalogo.tex}, «A CristalChain --- um livro-razão cuja cifra
conserva»).
A peça técnica é uma linha. O nome formal e a marca separam-se:

\begin{observacao}[Cadeia de Cristais (CristalChain)]\label{cripto:obs:nome}
\caixa{%
\textbf{Cadeia de Cristais}
$\;=\;$
nome formal da construção
\qquad
\textbf{CristalChain}
$\;=\;$
marca / repo / catálogo%
}
No primeiro aparecimento: \textbf{Cadeia de Cristais (CristalChain)}.
No texto técnico, \textbf{Cadeia de Cristais}.
Em produto, implementação, repositório ou marca, \textbf{CristalChain}.
O catálogo e a wiki já usam \code{CristalChain}; este paper não os
reescreve.
\end{observacao}

\begin{definicao}[livro-razão]\label{cripto:def:livro-razao}
Um \textbf{livro-razão} neste paper é uma sequência finita de
registos
\[
R_0\to R_1\to R_2\to\cdots\to R_m
\]
em que cada $R_i$ guarda o estado da órbita e o trabalho liquidado.
O sucessor obtém-se por um passo da cifra, \textbf{não} por consenso
entre pares.
Operacionalmente, $R_i$ é uma liquidação
(\code{tests/smartcontract.c}: estrutura \texttt{Liquidacao} ---
\texttt{liquidado}, régua, agente, \texttt{trabalho}).
A primeira liquidação acrescenta trabalho ao livro; as seguintes
são idempotentes (§S4: só a primeira conta).
Isto \textbf{não} é um bloco de rede nem um saldo económico.
\end{definicao}

\begin{observacao}[realização do livro-razão por identidade]\label{cripto:obs:livro-identidade}
O motor realiza o livro (Def.~\ref{cripto:def:livro-razao})
no \code{id} da identidade
(\code{app/src/banco\_identidade\_u.js}):
a chave vence a sessão; o livro pede dono.
\[
\boxed{\;
\begin{array}{c}
\mathrm{id}\text{ identifica o livro;}\\
R_i\text{ registra a liquidação;}\\
B\text{ fornece a cifra;}\\
\det(B^{k})=+1.
\end{array}
\;}
\]
Medido em \code{app/src/banco\_cristalchain\_u.js}
e \code{tests/cristalchain\_u.js} ($20/0$).
O nodo $\mathcal{U}$ usa \code (já existente);
não se inventa kind.
Idempotência da segunda liquidação
(\code{acrescentou}$=0$) \textbf{não} é dupla despesa ---
dupla despesa é \texttt{N/A}
(Obs.~\ref{cripto:obs:nao-blockchain}).
$\mathrm{sha256}$ produz o \code{id}; \textbf{não} é a cifra.
Pisano $\pi(q)$ conta a órbita do ouro na régua $\{B=1,\,C=-1\}$;
\textbf{não} é a companheira $B$ do ouro branco
(Obs.~\ref{cripto:obs:passo}).
\end{observacao}

\begin{definicao}[Cadeia de Cristais]\label{cripto:def:cadeia}
Uma \textbf{Cadeia de Cristais} (\textbf{CristalChain}) é um
livro-razão (Def.~\ref{cripto:def:livro-razao}) cuja cifra conserva
o invariante do passo (Def.~\ref{cripto:def:cifra-conserva}):
\[
\boxed{\;
\text{Cadeia de Cristais}
=
\text{livro-razão}
+
\text{cifra conservativa}
+
\det=+1
\;}
\]
A cifra \textbf{é forçada}, não escolhida.
Uma corrente cifrada por $\det=-1$ inverte o invariante a cada passo.
Uma corrente cifrada por $\det=+1$ conserva-o.
\textbf{Conservar} aqui é $\det(B^{k})=+1$ para todo $k$
(Teor.~\ref{cripto:thm:det}) --- o sinal do determinante sob
composição. \textbf{Não} é conservação de valor económico
(«nada se cria, nada se perde» no catálogo lê-se nesta restrição,
não como lastro).
\end{definicao}

\begin{observacao}[cifra $\to$ registo $\to$ cadeia]\label{cripto:obs:passo}
O passo algébrico da cifra do ouro branco é a companheira $B$:
\[
B=\begin{pmatrix}4&-1\\1&0\end{pmatrix},\qquad
\begin{pmatrix}a_{k+1}\\b_{k+1}\end{pmatrix}
=B\begin{pmatrix}a_{k}\\b_{k}\end{pmatrix},\qquad
\det(B^{k})=+1.
\]
Medido em \code{tests/ouro\_branco.c} §B2--§B3 (potências do $\det$).
A liquidação medida em \code{tests/smartcontract.c} dá o passo por
produto na régua: $\sigma=(0,1)$, \texttt{prod}$(r,a,b,0,1)$;
o contrato \textbf{regista} os passos como \texttt{trabalho}.
A paragem §S6 corre na régua do ouro $\{B=1,\,C=-1\}$
(Fibonacci $\bmod q$), \textbf{não} na companheira do ouro branco.
\caixa{%
passo da cifra ($B^{k}$, $\det=+1$)
$\;\neq\;$
passo da liquidação ($\sigma$, $\pi(q)$)%
}
Não se identificam as duas matrizes. Amarram-se os labels:
Def.~\ref{cripto:def:cifra-conserva} $\leftrightarrow$
\code{ouro\_branco.c};
Def.~\ref{cripto:def:livro-razao} $\leftrightarrow$
\code{smartcontract.c} \texttt{trabalho}.
\end{observacao}

\noindent
O enredo chama-lhe «ouro externo».
Isso é realização narrativa.
O objecto técnico é a Def.~\ref{cripto:def:cadeia}.
\textbf{Não} é analogia blockchain --- o próprio catálogo já o cortou
(Obs.~\ref{cripto:obs:nao-blockchain}).
O sítio canónico deste enunciado é \emph{aqui}
(\code{cripto:def:cadeia}); o catálogo e os medidores
\textbf{citam-se}, não se recopiam.

\begin{observacao}[dois objectos que partilham a palavra «cristal»]\label{cripto:obs:dois-cristais}
\caixa{%
Cadeia de Cristais (este paper, $\det=+1$)
$\;\neq\;$
corpus cristal (\code{cristal/cristal.jsonl}, volta por endereço)%
}
O corpus tem medidores próprios
(\code{tests/cristal\_volta.js}, \code{tests/cristal\_energia.js}):
reconstrução $R=0$ pelo \emph{id}, energia $E=\sum x^{2}$ conservada,
Parseval no anel $\mathbb{Z}/65537\mathbb{Z}$.
Isso é a volta do jornal, \textbf{não} a cifra do livro-razão.
Fundir os dois é homogeneizar homónimos.
Estatuto do elo «são o mesmo objecto»:
\texttt{nao localizada}
(§\ref{cripto:nao-localizada}).
\end{observacao}

% ═══════════════════════════════════════════════════════════════════════════
\section{Mapa do repositório --- citar, não recopiar}\label{cripto:mapa}

O que o repo já tem, por ofício. Cada linha é um objecto;
I0 proíbe colá-los.

\begin{center}\small
\begin{tabular}{@{}p{0.34\textwidth}p{0.58\textwidth}@{}}
\toprule
sítio & o que faz (e o que \textbf{não} é) \\
\midrule
\code{catalogo.tex} § CristalChain
  & enunciado-fonte: livro-razão cuja cifra conserva;
    \textbf{não} blockchain; fronteira dita \\
\code{tests/ouro\_branco.c}
  & $\det=\pm1$, norma, dobras $5,8,12$; cifra da Def.~\ref{cripto:def:cifra-conserva} \\
\code{tests/smartcontract.c}
  & paragem $\pi(q)$, agentes por $\Delta$, §S6 (teste $\neq$ prova) \\
\code{app/src/banco\_cristalchain\_u.js}
  & livro no disco GKBANCO; dono $=$ \code{id} da identidade;
    cinco contratos liquidam-se; \textbf{não} blockchain \\
\code{app/src/banco\_identidade\_u.js}
  & chave $>$ sessão; \code{id}$=$\code{hex}(banda);
    selector $\mathrm{sha256}$, \textbf{não} a cifra \\
\code{tests/cristalchain\_u.js}
  & medidor do motor + identidade \\
\code{tests/trialidade.c}
  & $S_3$ dos eixos; sobe a assinatura, \textbf{não} Duo \\
\code{tests/trialidade\_linguas.js}
  & $T^3$ das línguas \texttt{nao localizada} \\
\code{tests/familia\_real.c}
  & gato $A_m$, $\det=-1$, volta exacta; \textbf{não} sigilo (linear) \\
\code{tests/cifra\_entrada.c}, \code{lib/cifra.h}
  & cifra da ISA / fração contínua; \textbf{não} a cadeia \\
\code{tests/cifra\_prova.c}
  & cifra da demonstração (cadeia de passos $\neq$ cobertura) \\
\code{tests/cristalino.c}
  & esquilo $\det=+1$, $D<0$, ordem $4/6$;
    \textbf{não} ouro branco ($d=3$, disc~$12$) \\
\code{tests/cristal\_volta.js},
\code{tests/cristal\_energia.js}
  & volta do jornal \code{cristal.jsonl}; Parseval; objecto distinto \\
\code{tests/cristal\_avanco.js},
\code{tests/cristal\_curadoria.js},
\code{tests/cristal\_adversarial.js}
  & lotes / fusões / portão do corpus; \textbf{não} o livro-razão \\
\code{tests/cristal\_front.js}, \code{app/src/cristal.js}
  & chip X da UI (proveniência); \textbf{não} a cifra \\
\code{tests/dtcn.c}
  & $\star_s$, imposto algébrico; Paper~H; \textbf{não} livro-razão \\
\code{app/src/banco\_lei\_local\_u.js}
  & candidatos $\mathcal{L}_S$; $\mathrm{sha256}$ selector; sem promoção \\
\code{banco/bai.sql}
  & analogia «como a cristalchain»; guarda cifra do \emph{gato};
    \textbf{não} mede $\det=+1$ \\
\code{banco/sql.c} \code{OP\_MARTELO}
  & Hashcash no banco: $H(\mathrm{header})<T$ (80+32 átomos);
    \textbf{não} a Cadeia de Cristais (Obs.~\ref{cripto:obs:martelo}) \\
\code{cristal/cristal\_*.tex}, \code{cristal.jsonl}
  & wiki: nome, joalheria, GoldChain$\to$classe; papers prometidos ausentes \\
\code{enredo.tex}, \code{curriculo/}
  & ouro externo / proposta aberta; narrativa, não medidor \\
\code{app/src/manifesto.json}
  & dádiva $\sigma_1=\varphi$; enredo, não teorema \\
\code{hiper/.../paper\_K\_cripto.tex}
  & plano NIST + blockchain Patria; \textbf{outro} objecto
    (Obs.~\ref{cripto:obs:paper-k}) \\
\code{hiper/.../U4\_cripto*}
  & sandbox / células de literatura; sem incidência na Def.~\ref{cripto:def:cadeia} \\
\code{fisica.tex}, \code{corpo\_universal.tex},
\code{banco/} (código C)
  & zero ocorrências do nome; a ponte é o motor JS \\
\code{papers/cristalchain.tex},
\code{papers/cifra\_de\_cristal.tex}
  & referenciados na wiki; \textbf{ausentes} --- o canónico é este ficheiro \\
\bottomrule
\end{tabular}
\end{center}

\begin{observacao}[cristalino $\neq$ ouro branco]\label{cripto:obs:cristalino}
\code{tests/cristalino.c} mede o esquilo $\det=+1$ em $\mathbb{Q}(\sqrt D)$, $D<0$
(Gauss, Eisenstein; ordem $4$ ou $6$).
O ouro branco é $x^{2}=4x-1$ em $\mathbb{Z}[\sqrt3]$, disc~$12>0$.
Partilham o sinal do determinante; \textbf{não} são o mesmo corpo.
Elo «cristalino $=$ cifra da cadeia»: \texttt{nao localizada}.
\end{observacao}

\begin{observacao}[a família real não é sigilo]\label{cripto:obs:familia-real}
\code{tests/familia\_real.c}: $A_m$ é bijecção inteira e homomorfismo aditivo.
Isso dá reversibilidade e soma em cofre; \textbf{não} dá sigilo
(transformação linear cede a pares claro/cifrado).
O sigilo, se existir, vem de outro sítio --- e esse sítio,
como elo da Def.~\ref{cripto:def:cadeia}, está
\texttt{nao localizada}.
\end{observacao}

\begin{observacao}[Paper~K e U4 não são esta cadeia]\label{cripto:obs:paper-k}
\code{hiper/aposentados/teoria/papers/paper\_K\_cripto.tex}
é um plano (Kyber/Dilithium, $\star_p$, blockchain Substrate).
Não menciona CristalChain.
U4 é literatura viva no sandbox.
I0: plano $\neq$ Def.~\ref{cripto:def:cadeia};
blockchain genérica já é \texttt{N/A}
(Obs.~\ref{cripto:obs:nao-blockchain}).
\end{observacao}

% ═══════════════════════════════════════════════════════════════════════════
\section{A cifra que conserva}\label{cripto:cifra}

\begin{definicao}[cifra que conserva]\label{cripto:def:cifra-conserva}
No suporte $\mathbb{Z}[\sqrt d]$ a norma é $N(a+b\sqrt d)=a^{2}-d\,b^{2}$.
A matriz do passo é $2\times2$.
\[
\boxed{\;
\underbrace{x^{2}=nx+1}_{\text{família metálica},\ \det=-1}
\qquad\text{contra}\qquad
\underbrace{x^{2}=4x-1}_{\text{ouro branco},\ \det=+1}
\;}
\]
Conservar $=$ $\det=+1$ e $N=-1$ vazio
(\textbf{não} um valor económico; o sinal sob $B^{k}$).
Anti-conservar $=$ $\det=-1$ (o gato inverte a cada passo ímpar).
\end{definicao}

\begin{teorema}[o determinante decide o lado]\label{cripto:thm:det}
Para $n=1,2,3,4$ a matriz $A=\bigl(\begin{smallmatrix}n&1\\1&0\end{smallmatrix}\bigr)$
tem $\det=-1$.
A matriz $B=\bigl(\begin{smallmatrix}4&-1\\1&0\end{smallmatrix}\bigr)$
do ouro branco tem $\det=+1$.
$\det(A^{k})$ alterna com $k$; $\det(B^{k})=+1$ para todo $k$.
Não existe $n\in[-20,20]$ que ponha $2+\sqrt3$ na forma $x^{2}=nx+1$.
\end{teorema}

\begin{proof}
Conta na matriz, em $\mathbb{Z}$.
$(2+\sqrt3)^{2}=4(2+\sqrt3)-1$ exacto em $\mathbb{Z}[\sqrt3]$
(\code{tests/ouro\_branco.c} §B1--§B3).
\end{proof}

\begin{teorema}[a norma confirma por caminho independente]\label{cripto:thm:norma}
Nos unitários, $N$ só toma os três valores do trial.
O zero tem sempre exactamente uma solução (ponto fixo).
Quais dos outros dois $N$ atinge \textbf{é} a assinatura:
ouro ($d=5$) e prata ($d=2$) atingem $N=-1$ e $N=+1$;
ouro branco ($d=3$) atinge só $N=+1$.
\end{teorema}

\begin{proof}
Varredura em $a,b\in[-60,60]$
(\code{tests/ouro\_branco.c} §B6).
Duas contas que não se falam --- determinante e norma --- e o mesmo
veredito: o ouro branco não tem o lado que anti-conserva.
\end{proof}

\medido{$\det=-1$ em toda a família metálica e $+1$ no ouro branco;
$\det(A^{k})$ alterna na família e nunca no ouro branco;
$(2+\sqrt3)^{2}=4(2+\sqrt3)-1$ e zero soluções $x^{2}=nx+1$;
discriminantes $5$, $8$, $12$ (as dobras);
entre $n=1\ldots9$ um só $x^{2}-nx+1$ dá discriminante $12$
(\code{tests/ouro\_branco.c}).}

\noindent
As dobras $5$, $8$, $12$ são os discriminantes $b^{2}-4ac$,
calculados, não tabelados.
\textbf{Não} são Duo.
\textbf{Não} são $\operatorname{Star}(\mathcal{U})$.
Hexal $6=\mathrm{lcm}(2,3)$ é rótulo noutro suporte
(\code{univ:def:lcm}), não entra nesta ficha.

% ═══════════════════════════════════════════════════════════════════════════
\section{O que a cadeia não é}\label{cripto:fronteira}

\begin{observacao}[não é blockchain]\label{cripto:obs:nao-blockchain}
O catálogo e o medidor já o demonstraram:
não há rede, nós, nem consenso entre pares;
isto \textbf{não} resolve dupla despesa.
\caixa{%
Cadeia de Cristais
$\;\neq\;$
blockchain
$\;\neq\;$
GoldChain / HTLC / PoW
(\code{cristal/cristal\_computacao.tex})%
}
GoldChain nomeia $\sigma_{1}$ (um metal).
A wiki promoveu o nome à classe cristal
(\code{cristalchain} especializa \code{goldchain}).
Isso é nomenclatura do jornal, \textbf{não} prova de rede.
Estatuto de «rede que não racha / bloco minerado a frio»
como realização da Def.~\ref{cripto:def:cadeia}:
\texttt{nao localizada}.
PoW desmonta-se à parte (Obs.~\ref{cripto:obs:pow}):
Hashcash pode existir sem cadeia Nakamoto,
e o \code{MARTELO} do banco não é este livro.
\end{observacao}

\begin{proposicao}[a fronteira faz parte do resultado]\label{cripto:prop:fronteira}
Está estabelecido: a cifra (Teor.~\ref{cripto:thm:det},~\ref{cripto:thm:norma})
e a paragem (Teor.~\ref{cripto:thm:paragem}).
O resto --- consenso, dupla despesa, computação geral, oráculo do mundo ---
é trabalho por fazer ou recusa.
A linguagem do contrato exprime corpos de grau~$2$, não Turing geral.
O que se verifica a cada corrida é teste, não prova formal
(\code{tests/smartcontract.c} §S6).
\end{proposicao}

% ═══════════════════════════════════════════════════════════════════════════
\section{Paragem, reentrância, oráculo interno}\label{cripto:contrato}

O contrato \textbf{não se assina: liquida-se}.
Fornece-se metade; o código deriva o resto e executa
(\code{tests/smartcontract.c}).
O agente não vem da entrada: $\Delta$ despacha
gira / estica / limite.

\begin{teorema}[a paragem é da álgebra]\label{cripto:thm:paragem}
Num corpo finito $\bmod q$, a órbita do ouro fecha e não passa o tecto $q^{2}$.
O número de passos é o período de Pisano $\pi(q)$, obtido pela
recorrência de Fibonacci $\bmod q$ \textbf{sem tocar na órbita}.
\end{teorema}

\begin{proof}
Duas contas independentes, $q\in\{3,4,5,7,11,12\}$:
órbita por produto na régua $\{B=1,C=-1\}$ contra Fibonacci $\bmod q$.
Passos $=8,6,20,16,10,24=\pi(q)$.
Mutação $\bmod(q{+}1)$ faz a asserção falhar
(\code{tests/smartcontract.c} §S6).
\end{proof}

\medido{órbita fecha para $q\in\{3,4,5,7,11,12\}$ sob $q^{2}$;
passos batem com $\pi(q)$ por Fibonacci; o agente é função de $\Delta$,
não da entrada (quatro entradas do ouro, um só agente);
liquidar duas vezes é idempotente
(\code{tests/smartcontract.c}).}

\noindent
A reentrância não existe \emph{neste} suporte: a entrada não escolhe o
agente.
O oráculo deixa de ser preciso \emph{onde a metade que falta se deriva}
da metade fornecida, resíduo~$0$.
\textbf{Isto não resolve o problema do oráculo em geral}:
dados do mundo continuam a precisar de fonte.
Dizer o contrário seria vender o que não há.

\begin{observacao}[o trial sobe a cadeia, sem acrescentar]\label{cripto:obs:trial}
A assinatura tem três entradas porque o trial tem três estados
($p$ eixos em $+1$, $q$ em $-1$, $r$ em $0$).
Cada eixo é um trial.
\[
\boxed{\;
3\text{ autoduais}
\to
2^{3}=8\text{ biduais}
\to
8\times4=32\text{ órbitas}
\to
16{+}16\text{ pela paridade}
\;}
\]
Medido em \code{tests/trialidade.c}.
A trialidade é $S_{3}$ nos três eixos.
\textbf{Não} é $T^{3}$ das línguas
(\code{tests/trialidade\_linguas.js}:
trialidade tripla \texttt{nao localizada}).
Não se funde com Duo.
\end{observacao}

% ═══════════════════════════════════════════════════════════════════════════
\section{PoW --- as quatro peças não se fundem}\label{cripto:pow}

O PoW clássico é \textbf{uma} desigualdade, não quatro coisas coladas.
Nakamoto (2008), Hashcash (Back, 1997):
\[
\boxed{\;
H(\mathrm{header})<T
\;}
\]
O \emph{header} inclui o nonce; $H=\mathrm{SHA256}^{2}$;
$T$ é o alvo, função da dificuldade $d$.
A prova é o par $(\mathrm{header},n)$ que satisfaz a desigualdade.
A busca é sempre \emph{forward} (avalia-se $H$, não se inverte).
A verificação é uma avaliação.
O custo médio é $\mathbb{E}[\#\text{tentativas}]\propto T^{-1}$.

Isto \textbf{não} é a Def.~\ref{cripto:def:cadeia}.
Quatro ofícios, e nenhum implica o seguinte:

\begin{observacao}[quatro ofícios, não se fundem]\label{cripto:obs:pow}
\caixa{%
trabalho (busca de $n$)
$\;\neq\;$
verificação (um $H$)
$\;\neq\;$
selecção de cadeia (a mais pesada)
$\;\neq\;$
consenso (pares convergem)%
}
Uma cadeia Nakamoto \emph{usa} as quatro.
Hashcash usa só as duas primeiras (spam de correio, sem rede).
Consenso e selecção de cadeia já estão \texttt{N/A}
na Def.~\ref{cripto:def:cadeia}
(Obs.~\ref{cripto:obs:nao-blockchain}; §S6).
As duas primeiras não se herdam por homónimo.
\end{observacao}

\begin{observacao}[I0 --- três «trabalhos», três hashes]\label{cripto:obs:pow-i0}
\caixa{%
\begin{tabular}{@{}lll@{}}
$H(\mathrm{header})<T$
  & Hashcash / Nakamoto
  & busca; testemunha $(h,n)$ \\
$\mathrm{sha256}(\mathit{pub})\to\mathrm{id}$
  & selector da identidade
  & \textbf{não} a cifra; \textbf{não} o alvo \\
\texttt{trabalho} em \code{smartcontract.c}
  & passos da órbita
  & \textbf{não} hashcash \\
\end{tabular}%
}
$\mathrm{kind}=\mathrm{orbita}$ (já existente)
$\neq$ cadeia Nakamoto (\texttt{prevhash}, altura, a mais pesada).
$\pi(q)$ é o período da paragem (Teor.~\ref{cripto:thm:paragem}),
\textbf{não} o alvo $T$.
$\det(B^{k})=+1$ é identidade de conservação,
\textbf{não} um limiar a bater.
O \emph{gas} da EVM é cerca por fora da paragem Turing ---
o contrato já o recusa (§S6);
\textbf{não} é dificuldade de mineração.
\end{observacao}

\begin{observacao}[cinco peças na Cadeia de Cristais]\label{cripto:obs:pow-pecas}
Cada linha é um ofício.
O que a cadeia \textbf{já possui} não se homogeneíza com a coluna da esquerda.

\begin{center}\small
\begin{tabular}{@{}p{0.22\textwidth}p{0.14\textwidth}p{0.54\textwidth}@{}}
\toprule
peça (PoW clássico) & estatuto & evidência neste suporte \\
\midrule
prova de trabalho
  (custo de busca verificável)
  & \texttt{N/A}
  & \texttt{trabalho} $=$ passos até a órbita fechar
    (\code{agente\_corre}); mesma entrada, mesmo número
    (§S2; §S6, 200 corridas).
    Não há nonce, não há faixa, não há tentativa rejeitada.
    \code{liquida} é função pura dos termos. \\
dificuldade
  (parâmetro que encarece achar)
  & \texttt{N/A}
  & $\pi(q)$ é período, obtido por Fibonacci $\bmod q$
    \textbf{sem} tocar na órbita --- número com nome, não botão $T$.
    $q$ é o módulo do corpo.
    $\det=+1$ vale para todo $k$, exacto, sem limiar.
    Não se inventa $d$ por analogia. \\
verificação barata
  (verificar $\ll$ produzir)
  & \texttt{N/A}
  & o segundo caminho (Fibonacci) é a \textbf{mesma classe}
    $O(\pi(q))$ que a órbita: duas contas que podiam discordar
    (Teor.~\ref{cripto:thm:paragem}), não uma avaliação contra
    uma busca.
    Relêr o livro é ler um inteiro gravado, não conferir testemunha.
    Obs.~\ref{cripto:obs:pow-contas}. \\
selecção de cadeia
  (várias histórias $\to$ critério)
  & \texttt{N/A}
  & um livro por dono; o sucessor é liquidação
    (Def.~\ref{cripto:def:livro-razao}), não a cadeia mais pesada.
    Sem forks no motor
    (\code{tests/cristalchain\_u.js} §C7: donos distintos não se misturam).
    Não há histórias a escolher. \\
consenso
  (vários participantes convergem)
  & \texttt{N/A}
  & sem rede, sem pares
    (Obs.~\ref{cripto:obs:nao-blockchain}; §S6).
    Idempotência da liquidação $\neq$ dupla despesa. \\
\bottomrule
\end{tabular}
\end{center}

Nenhuma peça é \textbf{realizada} como PoW, nem \textbf{derivável}
de $\det=+1$ ou de $\pi(q)$ sem inventar o passo.
\end{observacao}

\begin{observacao}[duas contas independentes $\neq$ assimetria PoW]\label{cripto:obs:pow-contas}
A órbita $\bmod q$ e Fibonacci $\bmod q$ são dois caminhos
do \textbf{mesmo} número $\pi(q)$, ambos lineares no período.
É o padrão de medição já usado no determinante contra a norma
(Teor.~\ref{cripto:thm:det},~\ref{cripto:thm:norma}):
contas que não se falam e podem discordar.
Isso \textbf{não} é $\text{verificar}\ll\text{produzir}$.
A assimetria Hashcash pede uma via de sentido único e um limiar;
aqui a paragem é gaiola finita, e \textbf{toda} corrida fecha
(Teor.~\ref{cripto:thm:paragem}).
Elo algébrico «$\pi(q)$ realiza $H(\mathrm{header})<T$»:
\texttt{N/A} --- o suporte demonstra outro objecto.
\end{observacao}

\begin{observacao}[\code{MARTELO} $\neq$ livro-razão]\label{cripto:obs:martelo}
\code{banco/sql.c} \code{OP\_MARTELO} realiza Hashcash:
cabeçalho de 80~átomos, alvo de 32, $H=\mathrm{SHA256}^{2}$,
faixa de nonces, decisão na régua elíptica
(«a dificuldade é o raio da bola»).
Isso é PoW no banco ISA --- objecto \textbf{distinto}.
\code{app/src/banco\_cristalchain\_u.js} não o chama:
sem \code{S\_CAB}, sem nonce, sem $T$.
\caixa{%
\code{MARTELO} (Hashcash no banco)
$\;\neq\;$
\texttt{trabalho} da liquidação
$\;\neq\;$
Def.~\ref{cripto:def:cadeia}%
}
Elo «\code{MARTELO} $=$ cifra da cadeia»:
\texttt{nao localizada}
(o martelo existe; o pipe $\det=+1$ não o atravessa;
a wiki «bloco minerado a frio» já estava neste estatuto).
\end{observacao}

% ═══════════════════════════════════════════════════════════════════════════
\section{GLH, Parseval, lei local --- elos e cortes}\label{cripto:elos}

\subsection{GLH e $P=k/N$}
$\mathrm{I}_{4}=\mathrm{GLH}$ é o rectângulo das duas contagens
(\code{univ:def:glh-an}; \code{fis:thm:central}).
$P=k/N$ é Hurwitz nesse rectângulo
(\code{fis:lem:conserva}).
A cadeia \textbf{conta} órbitas (Pisano, Teor.~\ref{cripto:thm:paragem})
e \textbf{conserva} por $\det=+1$.
Contar $\neq$ conservar o sinal do determinante.
\[
\boxed{\;
P=k/N\ (\mathrm{GLH})
\neq
\det=+1\ (\text{cifra})
\neq
\lvert\det\rvert=1\ (\text{Lei~4 / volume})
\neq
p_i=\lvert\langle i\rvert\psi\rangle\rvert^{2}\ (\text{Born})
\;}
\]
Elo «a cadeia \emph{é} GLH»: \texttt{nao localizada}.
O que está realizado: as duas contas não se falam e não se fundem.

\subsection{Parseval do corpus}
\code{tests/cristal\_energia.js} mede
$E(\text{fonte})=E(\text{reconstrução})$ (inteiro exacto)
e Parseval $1\mathrm{D}/2\mathrm{D}/4\mathrm{D}$ no anel $p=65537$,
factor $256$, raiz errada quebra.
Isso é a volta do jornal cristal
(Obs.~\ref{cripto:obs:dois-cristais}),
\textbf{não} a cifra da Def.~\ref{cripto:def:cifra-conserva},
\textbf{não} Born,
\textbf{não} Lei~8 (não há Lei~8).
Elo «Parseval do corpus $=$ cifra da cadeia»:
\texttt{nao localizada}.

\subsection{Lei local e $\mathrm{sha256}$}
\code{banco\_lei\_local\_u.js} varre candidatos de
\code{univ:def:lei-local-canonica}.
$\mathrm{id}=\mathrm{hex}(\mathrm{sha256}(\mathrm{bytes}(\mathit{pub})))$
é o selector discreto da partição $B(\mathrm{id})$,
\textbf{não} o espectro $\mathcal{F}$.
A banda do canal é $\mathrm{sha256}$; $B_{\mathrm{cliente}}$ é outra
leitura (\code{fis:def:transf}).
Todos os candidatos saem \texttt{nao localizada};
sem Ficha~11; sem corpo novo.
\caixa{%
cifra da cadeia ($\det=+1$)
$\;\neq\;$
selector $\mathrm{sha256}$
$\;\neq\;$
$\mathcal{L}_{S}(\mathrm{id})$
$\;\neq\;$
$\mathcal{F}$ (\code{fis:thm:H})%
}
São o mesmo objecto? \texttt{nao localizada}
--- e o motor já recusa a promoção
(\code{banco\_lei\_local\_u.js} \code{RECUSAS};
\code{semPromocoes}).

\subsection{Cifra do gato (simétrica)}
A wiki (\code{cristal/cristal\_computacao.tex}) descreve
$A_{n}=\bigl(\begin{smallmatrix}n&1\\1&0\end{smallmatrix}\bigr)$,
$\det=-1$, bijecção $\bmod N$: cifrar $=A_{n}^{k}$,
decifrar $=(A_{n}^{-1})^{k}$.
Isso é a família metálica --- a que \textbf{anti-conserva}
(Teor.~\ref{cripto:thm:det}).
Pode ser cifra simétrica noutro ofício.
\textbf{Não} é a cifra da Cadeia de Cristais.
O paper prometido \code{papers/cifra\_de\_cristal.tex} não está no disco.
Dureza «idempotente $\Leftrightarrow$ factorar» (wiki, $143\to11\times13$):
sem medidor neste repo com o mesmo objecto que a Def.~\ref{cripto:def:cadeia}
--- \texttt{nao localizada} como elo da cadeia.

\subsection{DTC hipercomplexo}
\code{tests/dtcn.c} mede
$\lVert z\star_{s}w\rVert^{2}+(1-s^{2})\lVert a\times b\rVert^{2}
=\lVert z\rVert^{2}\lVert w\rVert^{2}$
(Paper~H, \code{hiper/aposentados/teoria/papers/paper\_H\_dtc\_hipercomplexo.tex}).
Conservação algébrica de motor, \textbf{não} livro-razão.
Elo com a cadeia: \texttt{nao localizada}.

% ═══════════════════════════════════════════════════════════════════════════
\section{O que ficou \texttt{nao localizada}}\label{cripto:nao-localizada}

\begin{center}\small
\begin{tabular}{@{}p{0.46\textwidth}p{0.46\textwidth}@{}}
\toprule
candidato de elo & porquê \\
\midrule
Cadeia de Cristais $=$ blockchain / consenso / dupla despesa
  & catálogo e §S6 \textbf{demonstram a ausência}: é \texttt{N/A},
    não busca falhada \\
Cadeia de Cristais $=$ PoW / Hashcash / $H(\mathrm{header})<T$
  & cinco peças \texttt{N/A} (Obs.~\ref{cripto:obs:pow-pecas});
    \texttt{trabalho} C $\neq$ work Bitcoin;
    $\mathrm{sha256}\to\mathrm{id}\neq$ alvo \\
$\pi(q)$ $=$ dificuldade $T$;
    $\det=+1$ $=$ limiar
  & período e conservação, não busca;
    analogia recusada (Obs.~\ref{cripto:obs:pow-i0}) \\
órbita $\parallel$ Fibonacci $=$ verificar $\ll$ produzir
  & mesma classe $O(\pi(q))$; medição, não assimetria
    (Obs.~\ref{cripto:obs:pow-contas}) \\
\code{MARTELO} $=$ cifra da cadeia
  & Hashcash no banco ISA; o motor da cadeia não o chama
    (Obs.~\ref{cripto:obs:martelo}) --- \texttt{nao localizada} \\
Cadeia de Cristais $=$ corpus cristal / Parseval do jornal
  & dois objectos; sem medidor que os identifique \\
Cadeia de Cristais $=$ $\mathcal{L}_{S}(\mathrm{id})$ / $\mathrm{sha256}$
  & selector $\neq$ cifra; candidatos recusados no motor \\
Cadeia de Cristais $=$ cifra do gato ($\det=-1$)
  & o gato anti-conserva; a cadeia pede $\det=+1$ \\
Cadeia de Cristais $=$ GLH / $P=k/N$
  & contar o rectângulo $\neq$ conservar o sinal \\
Cadeia de Cristais $=$ Born / palco quântico
  & $P=k/N\neq\lvert\psi_{i}\rvert^{2}$; Born não fecha $\mathcal{U}$ \\
Cadeia de Cristais $=$ DTC / $\star_{s}$
  & Paper~H é torque; sem incidência no livro-razão \\
Cadeia de Cristais $=$ Duo / $A_0$ / Lei~5 / Alonzo
  & cinco quadráticas (Obs.~\ref{cripto:obs:cinco-quadraticas}) \\
\code{papers/cifra\_de\_cristal.tex},
\code{papers/cristalchain.tex}
  & referenciados na wiki; \textbf{ausentes} no disco \\
rede «que não racha», bloco minerado a frio
    como realização da Def.~\ref{cripto:def:cadeia}
  & prosa da wiki; sem medidor da conservação $\det=+1$ nesse pipe \\
cristalino ($D<0$) $=$ ouro branco
  & partilham $\det=+1$; corpos distintos
    (Obs.~\ref{cripto:obs:cristalino}) \\
família real $A_m$ $=$ sigilo da cadeia
  & \code{familia\_real.c} mede o contrário: linear, sem sigilo \\
Paper~K / U4 $=$ Cadeia de Cristais
  & plano NIST/blockchain; sem o nome, sem o $\det=+1$ \\
BAI \code{banco/bai.sql} $=$ a cadeia
  & analogia + cifra do gato; sem medidor do ouro branco \\
\bottomrule
\end{tabular}
\end{center}

\noindent
\texttt{nao localizada} $=$ procurou-se; a evidência não apareceu.
\texttt{N/A} $=$ demonstrou-se que o suporte não admite o passo
(aqui: blockchain / consenso / dupla despesa / as cinco peças do PoW,
já cortados).
Não se promove coincidência de letras.

% ═══════════════════════════════════════════════════════════════════════════
\section{Medidores}\label{cripto:medidores}

\medido{Cifra e lado:
\code{tests/ouro\_branco.c} §B1--§B6.
Paragem, agentes, fronteira:
\code{tests/smartcontract.c} §S1--§S6.
Motor + identidade:
\code{app/src/banco\_cristalchain\_u.js};
\code{tests/cristalchain\_u.js}
($20/0$; dono $=$ \code{id}; chave $>$ sessão;
\code; \texttt{N/A} blockchain e PoW;
Obs.~\ref{cripto:obs:livro-identidade}).
PoW desmontado (Hashcash $\neq$ liquidação $\neq$ \code{MARTELO}):
Obs.~\ref{cripto:obs:pow}--\ref{cripto:obs:martelo}.
Trial da assinatura:
\code{tests/trialidade.c}
(e \code{tests/trialidade\_linguas.js}: $T^{3}$ não é a def).
Volta do corpus (objecto distinto):
\code{tests/cristal\_volta.js},
\code{tests/cristal\_energia.js}.
Lei local / $\mathrm{sha256}$ (objecto distinto):
\code{app/src/banco\_lei\_local\_u.js};
\code{univ:def:lei-local-canonica}.
DTC (objecto distinto):
\code{tests/dtcn.c}.
Cristalino / família real (objectos distintos):
\code{tests/cristalino.c}; \code{tests/familia\_real.c}.
Mapa: §\ref{cripto:mapa}.
Fecho de $\mathcal{U}$ (não reabrir):
\code{univ:thm:composicao-fechos};
\code{tests/grafo\_fecho\_u.js}.}

\colunas{Teor.~\ref{cripto:thm:det},~\ref{cripto:thm:norma},~\ref{cripto:thm:paragem};
Obs.~\ref{cripto:obs:i0},~\ref{cripto:obs:cinco-quadraticas},~\ref{cripto:obs:conservacoes},
~\ref{cripto:obs:nome},~\ref{cripto:obs:livro-identidade},~\ref{cripto:obs:passo},
~\ref{cripto:obs:pow},~\ref{cripto:obs:pow-i0},~\ref{cripto:obs:pow-pecas},
~\ref{cripto:obs:pow-contas},~\ref{cripto:obs:martelo};
Def.~\ref{cripto:def:livro-razao},~\ref{cripto:def:cadeia},~\ref{cripto:def:cifra-conserva}}
        {ouro branco $\det=+1$; órbita $\bmod q=\pi(q)$;
livro-razão por identidade; sem consenso, sem PoW e sem saldo económico}
        {\code{ouro\_branco.c}, \code{smartcontract.c}, \code{cristalchain\_u.js}}

\end{document}
